\documentclass[openany,fontset=fandol]{ctexbook}
\usepackage{amsmath,amssymb,amsthm}
\usepackage[unicode,hidelinks]{hyperref}
\setlength{\emergencystretch}{2em}
\raggedbottom
\newcommand{\OO}{\mathcal O}
\newcommand{\FF}{\mathcal F}
\newcommand{\GG}{\mathcal G}
\newcommand{\PP}{\mathbf P}
\DeclareMathOperator{\Spec}{Spec}
\DeclareMathOperator{\coprof}{coprof}
\newcommand{\tagitem}[2]{\par\medskip\noindent\textbf{#1\quad #2}\enspace}
\begin{document}
\title{代数几何基础第五卷预备笔记\\法文覆盖审计：泛截面的不可约性与连通性}
\author{第91批\quad 法、英、中文三方校订}
\date{2026年11月23日}
\maketitle

\section*{范围与接续}
逐页对读法文物理14页下半、15--16页、17页上端，以及英编稿物理12--13页中第3节，
覆盖定理3.1、引理3.1.1、推论3.2和3.3。法14--15的划销文字另记，不混入定稿。
第三辑合订本在 \texttt{v5\_03\_06.tex} 已有对应译文；本批为待合订替换稿，
不重复计新译页、不改1158页基线。下次从法文物理17页的20.4、4.1及英文物理13页第4节接续。

\section*{20.3\quad 泛超平面截面：几何不可约性与连通性}
记 $P=\PP^n_k$，$P^\vee$ 为超平面参数空间，$\eta$ 为其泛点，
$Y_\eta$ 为 $f:X\to P$ 的泛超平面截面。像的维数理解为其闭包的维数。

\tagitem{定理}{3.1（Bertini--Zariski）}
若 $\dim\overline{f(X)}\ge 2$，且 $X$ 几何不可约，则 $Y_\eta$ 在 $k(\eta)$ 上几何不可约。

\begin{proof}[原稿的函数域归约]
设 $K/k$ 为 $X$ 的函数域，$n=\dim P$。在 $P$ 中选取仿射坐标 $T_1,\ldots,T_n$
（所选无穷远超平面 $H^\infty$ 不包含整个 $f(X)$），并在 $P^\vee$ 中选取相应的仿射坐标 $S_1,\ldots,S_n$。
写 $t_i\in K$ 为 $T_i$ 沿 $f$ 在泛点处的像。于是 $Y_\eta$ 的函数域 $L$ 等同于整环
\[
 K[S_1,\ldots,S_n]\big/\bigl(\textstyle\sum_{i=1}^n t_i S_i-1\bigr)
\]
的分式域。因为 $\dim\overline{f(X)}>0$，并非所有 $t_i$ 都在 $k$ 上代数，
所以也不全为零；重排后可设 $t_n\ne 0$。由关系式消去 $S_n$，可得
\[
 L=K(S_1,\ldots,S_{n-1}),\qquad
 S_n=a_0+\sum_{i=1}^{n-1}a_iS_i,\qquad
 a_0=t_n^{-1},\quad a_i=-t_i t_n^{-1}.
\]
因此 $L/K$ 为纯超越扩张。另一方面，$k'=k(\eta)$ 等同于参数有理函数域
$k(S_1,\ldots,S_n)$，典范嵌入 $k'\hookrightarrow L$ 将参数 $S_i$ 送入该商环中的类，
再视为 $L$ 中的元素。\footnote{法15用同一个 $S_i$ 记映射两端；英编稿编者要求核对，
并在注11建议写成等价类，正文现明确区分。旧稿中“函数域”在 $X$ 可能非约化时取其泛点剩余域，
或先换为 $X_{\rm red}$；基变换后该替换不改变拓扑空间的不可约性判断。}
所以作为 $L$ 的子域，$k'$ 既可由所有这些 $S_i$ 生成，也可由前 $n-1$ 个
$S_i$ 和上述线性式生成。由 $t_n=1/a_0$ 和 $t_i=-a_i/a_0$，两组系数在 $k$ 上生成同一子域；
共同的超越次数就是 $\dim\overline{f(X)}$。
\footnote{法15下端划掉较弱的“$a_i$ 不全在 $k$ 上代数”，手写改为共同超越次数等于像的维数。
在有理函数表示中需要的正维性与使用下述引理所需的至少二维不是同一门槛。}
原稿建议把这一双有理描述至少作为（2.1）的推论，然后把（3.1）归结为下面的引理。
\end{proof}

\noindent\emph{译校说明。}
取 $X_{\rm red}$ 后，关联族是其上的射影丛，故整。它支配 $P^\vee$：
由（2.3），正维像的泛截面非空。这样，在 $X$ 泛点上写出的商环描述的是关联族函数域，
而支配性保证它包含 $k(\eta)$，且也是 $Y_\eta$ 的函数域。
故上述典范嵌入不是假设一个任意的线性商会保留所有参数的独立性。
几何不可约性可由此时的整概形的函数域是否为准素扩张判定。

\tagitem{引理}{3.1.1（Zariski）}
设 $k$ 为域，$K/k$ 为有限型域扩张，$m\ge0$ 为整数，$a_i\in K$（$0\le i\le m$），且
$\operatorname{trdeg}_k k(a_0,\ldots,a_m)\ge2$。记
\[
 L=K(S_1,\ldots,S_m),\qquad
 k'=k\left(S_1,\ldots,S_m,a_0+\sum_{i=1}^m a_iS_i\right)\subset L,
\]
其中 $S_i$ 为不定元。若 $K/k$ 为准素扩张，则 $L/k'$ 为准素扩张。
\footnote{法16手写把“不全代数”换成超越次数至少2，并划销一个根重数等式，不列入引理。
法本保留 $m\ge0$；若 $m=0$，单个 $a_0$ 不可能满足超越次数至少2，所以该分支无实例。}

原稿说，这一引理或近似版本在文献中颇常见，把证明出处留给收信人选择，自己不另找一份证明。

\noindent\emph{译校说明。}
按正式 EGA IV 第二部分（4.3.1），“准素”指扩张域中底域的可分代数闭包就是底域，
并非要求其中每个代数元都属于底域。例如特征 $p>0$ 的 $k=\mathbf F_p(u)$
在 $k(u^{1/p})$ 中没有新的可分代数元，却有不属于 $k$ 的纯不可分代数元。
英编稿注12对术语的猜测不作定义。英编稿所列查证线索为 Zariski《论文集》第1卷第174页、
第2卷第304页，Zariski--Samuel 的 Springer GTM 版第1卷第196页、第2卷第230页，
以及 Jouanolou《Bertini 定理及应用》定理3.6和第6节；这是该编辑的书目，若未逐项核原书，
不据此断言具体页上的命题逐字等价。参见另有明确证明的
\href{https://stacks.math.columbia.edu/tag/0G4C}{Stacks 37.32}（Bertini 的现代证明）
和 \href{https://stacks.math.columbia.edu/tag/00I2}{10.47}（几何不可约代数）。
本手稿对3.1只写到归约于引理，并未提供自足证明。

\tagitem{推论}{3.2}
设 $f$ 非分歧，或 $k$ 的特征为零；另设 $\dim\overline{f(X)}\ge2$。
若 $X$ 几何整，则 $Y_\eta$ 也几何整。\footnote{法16将打字的“浸入”改成手写“非分歧”，与英编稿一致。
应读成“（非分歧或特征零）且像至少二维”，不是由特征零单独推出结论。}
\begin{proof}[原稿论证]
几何整等价于几何不可约且可分（这里的可分即几何约化）。
\end{proof}
\noindent\emph{译注。}
几何不可约来自（3.1）。非分歧时，几何约化由（2.14）的相应结论逐点转移；
特征零时，$k(\eta)$ 为完美域，正则态射 $Y_\eta\to X$ 将源端的约化性传到截面，
而有限型完美域上的约化即几何约化。

\tagitem{推论}{3.3}
设 $k$ 代数闭，且对 $X$ 的每个不可约分支 $X_i$，均有
$\dim\overline{f(X_i)}\ge2$。设 $\sigma$ 为像闭包维数不大于零（包含空像）的各闭子集
所成的族。若对每个 $Z\in\sigma$，$X\setminus Z$ 均连通，则 $Y_\eta$ 在 $k(\eta)$ 上几何连通。
\footnote{法16打字条件为维数等于零，并保留像的闭包标记。本译将空像并入，避免将不相交的两分支
误作图中的邻接。在这里，若 $X$ 非空而不连通，即使只要求去掉像维数为零的闭集仍连通，
也会矛盾：取其中一侧一个闭点，去掉它仍保留该侧至少一个正维不可约分支的非空开部，
且另一侧不变。空 $X$ 的陈述按空集连通的通常约定单独处理。}
\begin{proof}
原稿援引拟与 Hartshorne 定理同列于第6段的分支链引理。
具体地，把有限多个不可约分支作为顶点，若某两个分支交集的像有正维闭包，则连一条边。
若图非连通，把顶点分成非空两组，对应并集分别为 $U_1,U_2$。
公共闭集 $Z=U_1\cap U_2$ 是所有跨组分支交集的有限并，故其像闭包维数不大于零。
每组至少一支没有被 $Z$ 覆盖，因为每支的像至少二维。
由此 $X\setminus Z$ 分成两个非空互不相交的闭集，矛盾。因此图连通。

各 $X_i$ 带约化结构；因 $k$ 代数闭且 $X_i$ 不可约，由（3.1），其泛截面几何不可约，
因而几何连通。相邻两个分支交集的泛截面依（2.3）非空，且在任意扩域后仍非空。
在 $k(\eta)$ 的代数闭包上沿这张连通图逐步取并，得到整个截面连通。
\end{proof}
\noindent 法17上端结束（3.3），其后手稿明确有20.4和4.1的编号。
英编稿第4节漏单列4.1，下一批以法文页为准保留编号。

\section*{划销文字与版面证据}
法14的（3.1）之后，有一段跨至页尾的划销旧论述：对比两个较早表述，提及 $X$ 与泛截面可分，
并准备从（2.1）展开函数域；可辨部分不是现行证明的一部分。
法15顶部另有已划销、题头编号为“推论2.8”的旧稿，内容为非嵌入分支及根重数、平坦性与逐层分解，
和现行（2.16）相关。保留两个摘页供复核，不把旧稿串进（3.1）。
法16引理内的划销根重数断言也另列，而非现行结论。
\end{document}
